\chapter{Лабораторная работа №1.3. Марковские цепи. Определение и построение}
\label{ch:chap3}

\section{Цель работы}
\label{sec:goal3}

    Освоить методы построения и анализа марковских цепей в MATLAB: построение матрицы переходов для сети с коммутацией пакетов, проверка матрицы на стохастичность, проверка цепи на эргодичность, визуализация графа сети. Научиться программировать функции проверки свойств марковских цепей.

\section{Задание к лабораторной работе}
\label{sec:task3}

    В соответствии с методическими указаниями \cite{andreev2024tmo} и вариантом №2 необходимо:

    \begin{enumerate}
        \item Построить сеть из $L = 15$ узлов в виде ориентированного графа:
        \begin{itemize}
            \item каждый узел должен иметь минимум 3 исходящих маршрута;
            \item каждый узел должен иметь хотя бы один входящий маршрут.
        \end{itemize}
        \item Построить матрицу переходов $T = \|p_{ij}\|$ размерности $L = 15$ с ненулевыми элементами, удовлетворяющими условиям стохастичности.
        \item Написать функции в MATLAB:
        \begin{itemize}
            \item \verb|stochastic(matrix)| — проверяет матрицу на стохастичность;
            \item \verb|ergodic(matrix, epsilon)| — проверяет цепь Маркова на эргодичность.
        \end{itemize}
        \item Проверить стохастичность и эргодичность цепи Маркова.
        \item Оформить полученные результаты: матрицу переходов, функции и результаты проверок.
    \end{enumerate}

\section{Ход выполнения работы}
\label{sec:work3}

\subsection{Построение матрицы переходов}
\label{subsec:matrix3}

    Для сети из $L = 15$ узлов выбрана кольцевая топология, при которой каждый узел $i$ имеет три исходящих маршрута: к узлам $(i \bmod L) + 1$, $((i+1) \bmod L) + 1$ и $((i+2) \bmod L) + 1$. Каждому маршруту присвоена вероятность $1/3$:

    \begin{equation}
    \label{eq:matrix}
        T_{i, j} = \begin{cases}
            \frac{1}{3}, & j \in \{\text{targets}(i)\}, \\
            0, & \text{иначе}.
        \end{cases}
    \end{equation}

    Матрица переходов $T$ имеет размерность $15 \times 15$. Сумма элементов каждой строки равна 1, что подтверждает стохастичность матрицы.

\subsection{Проверка стохастичности}
\label{subsec:stoch3}

    Функция \verb|stochastic| проверяет два условия:
    \begin{enumerate}
        \item все элементы матрицы лежат в диапазоне $[0; 1]$;
        \item сумма элементов каждой строки равна 1 с точностью $10^{-9}$.
    \end{enumerate}

    \begin{figure}[ht]
        \centering
        \includegraphics[width=0.9\textwidth]{lab1_3_matrix_output}
        \caption{Матрица переходов $T$ и результат проверки на стохастичность}
        \label{fig:lab1_3_matrix_output}
    \end{figure}

    На рисунке \ref{fig:lab1_3_matrix_output} представлен вывод матрицы переходов и сообщение о том, что матрица является стохастической.

\subsection{Проверка эргодичности}
\label{subsec:erg3}

    Функция \verb|ergodic| проверяет цепь Маркова на эргодичность: матрица возводится в степень $m = 200$, после чего проверяется, что все элементы $P^m$ положительны с точностью $\varepsilon = 10^{-5}$. Если условие выполнено, цепь эргодическая, и выводится стационарное распределение (первая строка $P^m$).

    \begin{figure}[ht]
        \centering
        \includegraphics[width=0.9\textwidth]{lab1_3_ergodic_output}
        \caption{Результат проверки цепи Маркова на эргодичность}
        \label{fig:lab1_3_ergodic_output}
    \end{figure}

    На рисунке \ref{fig:lab1_3_ergodic_output} представлен вывод: сообщение об эргодичности цепи и стационарное распределение.

\subsection{Визуализация графа сети}
\label{subsec:graph3}

    Для визуализации сети использована функция \verb|digraph| с раскладкой \verb|'force'|. Граф содержит 15 узлов и 45 направленных рёбер (по 3 исходящих из каждого узла).

    \begin{figure}[ht]
        \centering
        \includegraphics[width=0.9\textwidth]{lab1_3_graph}
        \caption{Граф сети из 15 узлов}
        \label{fig:lab1_3_graph}
    \end{figure}

    На рисунке \ref{fig:lab1_3_graph} представлен ориентированный граф сети с кольцевой топологией.

\section{Выводы}
\label{sec:concl3}

    В ходе лабораторной работы №1.3 построена матрица переходов для сети из 15 узлов с кольцевой топологией. Матрица проверена на стохастичность — все строки дают в сумме 1, все элементы лежат в диапазоне $[0; 1]$. Цепь Маркова проверена на эргодичность — матрица $P^{200}$ содержит только положительные элементы, что подтверждает эргодичность цепи и наличие стационарного распределения. Получены навыки работы с цепями Маркова, программирования функций проверки свойств и визуализации графов в MATLAB.

\endinput